\documentclass[oneside,a4paper,14pt]{extarticle}
\usepackage[T1,T2A]{fontenc}
\usepackage[left=2.5cm,right=1.5cm,top=2cm,bottom=2cm]{geometry}
\usepackage[utf8]{inputenc}
\usepackage[russian]{babel}
\usepackage{textcomp}
\usepackage{indentfirst}
\usepackage{graphicx}
\usepackage{mwe}
\usepackage{wrapfig}
\usepackage{caption}
\usepackage{amsmath}
\usepackage{amsfonts}
\usepackage{amsthm}
\usepackage{enumitem}
\usepackage{minted}
\usepackage[all]{xy}
\usepackage[breaklinks]{hyperref}
\usepackage{titlesec}
\titleformat{\section}
{\normalsize\bfseries}
{\thesection} {1em} {}
     \titleformat{subsection}
{\normalsize\bfseries}
 {\thesubsection} {1em} {}
    \titleformat{\subsubsection}
{\normalsize\bfseries}
{\thesubsection} {1em} {}

\newcommand{\solutionheading}[1] {
    \begin{flushleft}
    \hspace* {1cm}\textbf{#1}
    \end{flushleft}
 }

\renewcommand\baselinestretch{1.1}\normalsize
\setlength{\parindent}{1cm}
\usepackage{indentfirst}


\begin{document}
    \newpage\thispagestyle{empty}
    \begin{center}
        МИНИСТЕРСТВО НАУКИ И ВЫСШЕГО ОБРАЗОВАНИЯ\\
            РОССИЙСКОЙ ФЕДЕРАЦИИ
            ФЕДЕРАЛЬНОЕ ГОСУДАРСТВЕННОЕ БЮДЖЕТНОЕ\\
            ОБРАЗОВАТЕЛЬНОЕ
            УЧРЕЖДЕНИЕ ВЫСШЕГО ОБРАЗОВАНИЯ\\
             ГОСУДАРСТВЕННЫЙ УНИВЕРСИТЕТ»\\
            Институт математики и информационных систем\\
            Факультет автоматики и вычислительной техники\\
            Кафедра электронных вычислительных машин
    \end{center}
    \vspace{20mm}

    \begin{center}
        Отчёт по лабораторной работе №2\\
        по дисциплине\\
        <<Объектно-Ориентированное Программирование>>\\
    \end{center}
    \vspace{48mm}

    Выполнил студент гр. -2304-05-00 \hspace{7mm} \rule[-0,5mm]{16mm}{0.15mm}\,/С Т.Е./


    Проверил преподаватель \hfill  \rule[-0,5mm]{30mm}{0.15mm}\,/Ш Н.А./

    \vfill
    \begin{center}
        ов\\
        2026
    \end{center}
    \newpage

    \section*{Цель}
Изучить основы объектно-ориентированного программирования в языке C# на примере разработки приложения «Умный дом» с графическим интерфейсом и хранением данных в БД SQLite.

\section*{Задание}
На основании согласованной в первой лабораторной работе предметной
области разработать приложение с графическим интерфейсом пользователя,
используя классы и особенности C#.
Сведения об объектах предметной области необходимо хранить в БД.\\
1. Все требования 1ой работы в силе, каждый класс отдельно, базовый\\
абстрактный, наследование, доступ к полям классов через геттеры/сеттеры,
не напрямую.\\ 2. Есть графический интерфейс. По умолчанию WindowsForms.
Другие варианты обсуждаются, но это не консольное приложение.\\ 3. Для хранения данных используется БД. Данные хранятся не скопом в одной таблице.
Реализованы основные CRUD функции. Методы классов из 1ой работы тоже доступны.\\ 4. В отчете есть описание диаграммы классов(не забудьте про
классы формы, они тоже должны быть). Прописываем особенности
подключения к БД.


    \newpage

\section*{Решение}

\section*{Структура проекта:}

Выбранная предметная область: «Умный дом». Абстрактный класс \textbf{Device} и 4 дочерних класса — \textbf{Thermostat}, \textbf{Light}, \textbf{DoorLock} и \textbf{Camera}.
\begin{figure}[H]
    \centering
    \includegraphics[width=1\linewidth]{image1.png}
    \caption*{Рисунок 1 — Структура проекта}
    \label{fig:placeholder}
\end{figure}

\newpage

\begin{figure}[H]
    \centering
    \includegraphics[width=1\linewidth]{ClassDiagram.png}
    \caption*{Рисунок 2 — Диаграмма классов}
    \label{fig:placeholder}
\end{figure}


\newpage

\textbf{Абстрактный класс Device}\\

Базовый абстрактный класс для всех устройств умного дома.

Собственные поля:

\begin{verbatim}
string _name;          // Название устройства
string _manufacturer;  // Производитель
bool _isOn;            // Состояние (вкл/выкл)
DateTime _createdAt;   // Дата создания
\end{verbatim}

\textbf{Конструкторы:}

\begin{enumerate}
    \item Конструктор по умолчанию Device();

Нет входных параметров, значения полей \_name, \_manufacturer принимают значения по умолчанию, \_isOn = false, \_createdAt = текущее время.

\item Конструктор Device(string name, string manufacturer);

Поле \_name получает входное строковое значение name.\\
Поле \_manufacturer получает входное строковое значение manufacturer.\\
\_isOn = false, \_createdAt = текущее время.

\item Конструктор копирования Device(Device other);

Копирует значения всех полей из объекта other.

\end{enumerate}\\

\textbf{Геттеры/Сеттеры (свойства):}

\begin{verbatim}
string Name          — чтение/запись поля _name
string Manufacturer  — чтение/запись поля _manufacturer
bool IsOn            — чтение/запись поля _isOn
DateTime CreatedAt   — чтение/запись поля _createdAt
\end{verbatim}

\textbf{Методы:}

\begin{verbatim}
virtual void TurnOn()              — включить устройство
virtual void TurnOff()             — выключить устройство
virtual string GetStatus()         — получить статус
virtual string GetDeviceType()     — получить тип устройства
virtual string GetDetailedInfo()   — детальная информация
abstract string GetShortDescription()  — краткое описание
\end{verbatim}


\textbf{Дочерний класс Thermostat}\\

Наследник класса Device, представляет термостат — устройство управления температурой и влажностью.

Собственные поля:

\begin{verbatim}
double _currentTemperature;  // Текущая температура
double _targetTemperature;   // Целевая температура
string _mode;                // Режим (heat/cool/auto)
double _humidity;            // Влажность
\end{verbatim}

\textbf{Конструкторы:}\\

Thermostat() — по умолчанию, 20°C → 22°C, режим auto, влажность 50\%.\\

Thermostat(string name, string manufacturer, double targetTemp) — основной.\\

Thermostat(string name, string manufacturer, double targetTemp, double currentTemp, string mode, double humidity) — полный.\\

Thermostat(Thermostat other) — копирования.\\

\textbf{Геттеры/Сеттеры:}

\begin{verbatim}
double CurrentTemperature  — чтение/запись текущей температуры
double TargetTemperature   — чтение/запись целевой температуры
string Mode                — чтение/запись режима
double Humidity            — чтение/запись влажности
\end{verbatim}

\textbf{Собственные методы:}

\begin{verbatim}
void SetTemperature(double temp)  — установить целевую температуру
string GetClimateInfo()           — информация о микроклимате
\end{verbatim}

\textbf{Перекрытые методы (override):}

\begin{verbatim}
TurnOn()            — включение с запуском режима
GetDetailedInfo()   — детальная информация
GetDeviceType()     — возвращает "Thermostat"
GetShortDescription() — краткое описание
\end{verbatim}

Унаследованные методы: TurnOff(), GetStatus().

\newpage

\textbf{Дочерний класс Light}\\

Наследник класса Device, представляет умное освещение с регулировкой яркости и цвета.

Собственные поля:

\begin{verbatim}
int _brightness;      // Яркость (0-100)
string _color;        // Цвет
bool _autoMode;       // Авторежим
int _scheduledTime;   // Время включения по расписанию
\end{verbatim}

\textbf{Конструкторы:}\\

Light() — по умолчанию, яркость 50\%, цвет white, ручной режим.\\

Light(string name, string manufacturer, int brightness) — основной.\\

Light(string name, string manufacturer, int brightness, string color, bool autoMode, int scheduledTime) — полный.\\

Light(Light other) — копирования.\\

\textbf{Геттеры/Сеттеры:}

\begin{verbatim}
int Brightness       — чтение/запись яркости (0-100)
string Color         — чтение/запись цвета
bool AutoMode        — чтение/запись авторежима
int ScheduledTime    — чтение/запись времени расписания
\end{verbatim}

\textbf{Собственные методы:}

\begin{verbatim}
void SetBrightness(int level)  — установить яркость (0-100)
string GetLightInfo()          — информация об освещении
\end{verbatim}

\textbf{Перекрытые методы (override):}

\begin{verbatim}
TurnOn()            — включение с автоувеличением яркости
GetDetailedInfo()   — детальная информация
GetDeviceType()     — возвращает "Light"
GetShortDescription() — краткое описание
\end{verbatim}

Унаследованные методы: TurnOff(), GetStatus().

\newpage

\textbf{Дочерний класс DoorLock}\\

Наследник класса Device, представляет умный дверной замок с кодовым доступом.

Собственные поля:

\begin{verbatim}
bool _isLocked;           // Состояние (заперт/открыт)
string _accessCode;       // Код доступа
int _wrongAttempts;       // Неверные попытки
bool _autoLock;           // Автоблокировка
string _lastAccessTime;   // Время последнего доступа
\end{verbatim}

\textbf{Конструкторы:}\\

DoorLock() — по умолчанию, заперт, код "1234", автоблокировка включена.\\

DoorLock(string name, string manufacturer, string accessCode) — основной.\\

DoorLock(string name, string manufacturer, string accessCode, bool autoLock, bool isLocked) — полный.\\

DoorLock(DoorLock other) — копирования.\\

\textbf{Геттеры/Сеттеры:}

\begin{verbatim}
bool IsLocked          — чтение/запись состояния
string AccessCode      — чтение/запись кода доступа
int WrongAttempts      — чтение/запись неверных попыток
bool AutoLock          — чтение/запись автоблокировки
string LastAccessTime  — чтение/запись времени доступа
\end{verbatim}

\textbf{Собственные методы:}

\begin{verbatim}
bool TryUnlock(string code)  — попытка открыть по коду
string GetLockInfo()         — информация о замке
\end{verbatim}

\textbf{Перекрытые методы (override):}

\begin{verbatim}
TurnOn()            — включение с автоблокировкой
GetDetailedInfo()   — детальная информация
GetDeviceType()     — возвращает "DoorLock"
GetShortDescription() — краткое описание
\end{verbatim}

Унаследованные методы: TurnOff(), GetStatus().

\newpage

\textbf{Дочерний класс Camera}\\

Наследник класса Device, представляет камеру видеонаблюдения с детекцией движения.

Собственные поля:

\begin{verbatim}
int _resolution;         // Разрешение (720/1080/2160/3840p)
bool _isRecording;       // Состояние записи
bool _motionDetection;   // Детекция движения
int _storageUsed;        // Использовано памяти (%)
string _lastMotionTime;  // Время последнего движения
\end{verbatim}

\textbf{Конструкторы:}\\

Camera() — по умолчанию, 1080p, детекция включена, память 0\%.\\

Camera(string name, string manufacturer, int resolution) — основной.\\

Camera(string name, string manufacturer, int resolution, bool recording, bool motionDetect, int storageUsed) — полный.\\

Camera(Camera other) — копирования.\\

\textbf{Геттеры/Сеттеры:}

\begin{verbatim}
int Resolution          — чтение/запись разрешения
bool IsRecording        — чтение/запись состояния записи
bool MotionDetection    — чтение/запись детекции движения
int StorageUsed         — чтение/запись использованной памяти
string LastMotionTime   — чтение/запись времени движения
\end{verbatim}

\textbf{Собственные методы:}

\begin{verbatim}
string StartRecording()  — начать запись (с проверкой памяти)
string RecordMotion()    — зафиксировать движение
\end{verbatim}

\textbf{Перекрытые методы (override):}

\begin{verbatim}
TurnOn()            — включение с инициализацией режима
GetDetailedInfo()   — детальная информация
GetDeviceType()     — возвращает "Camera"
GetShortDescription() — краткое описание
\end{verbatim}

Унаследованные методы: TurnOff(), GetStatus().

\newpage

\textbf{DatabaseHelper.cs}\\

Статический класс для работы с базой данных SQLite. Все запросы выполняются напрямую через System.Data.SQLite (без Entity Framework).

\begin{enumerate}
\item \textbf{Initialize()} — Инициализация БД. Создаёт файл \texttt{smart\_home.db} и формирует таблицы через \texttt{CreateTables()}.

\item \textbf{CreateTables()} — Создание пяти таблиц: \texttt{Devices} (общие поля), \texttt{Thermostats}, \texttt{Lights}, \texttt{DoorLocks}, \texttt{Cameras}. Каждая дочерняя таблица связана с \texttt{Devices} через внешний ключ \texttt{DeviceId} с каскадным удалением.

\item \textbf{CRUD-операции:}
\begin{itemize}
    \item \texttt{InsertDevice(device)} — INSERT в Devices, возвращает Id.
    \item \texttt{InsertThermostat(thermostat, deviceId)} — INSERT специфичных данных термостата.
    \item \texttt{InsertLight(light, deviceId)} — INSERT данных освещения.
    \item \texttt{InsertDoorLock(doorLock, deviceId)} — INSERT данных замка.
    \item \texttt{InsertCamera(camera, deviceId)} — INSERT данных камеры.
    \item \texttt{GetAllDevices()} — SELECT * FROM Devices, возвращает DataTable.
    \item \texttt{GetDeviceById(id)} — SELECT с WHERE Id = ?, возвращает словарь.
    \item \texttt{GetThermostatByDeviceId(deviceId)} — SELECT из Thermostats.
    \item \texttt{UpdateDevice(device, id)} — UPDATE общих полей.
    \item \texttt{SetDeviceState(id, isOn)} — UPDATE IsOn (вкл/выкл).
    \item \texttt{DeleteDevice(id)} — DELETE из дочерней таблицы + Devices.
    \item \texttt{CreateDeviceObject(deviceId)} — создаёт объект C# нужного типа-наследника на основе данных из БД (через \texttt{switch} по \texttt{DeviceType}).
    \item \texttt{SaveDeviceObject(device, deviceId)} — сохраняет изменения объекта обратно в БД.
\end{itemize}

\end{enumerate}

\newpage
\begin{minted}{csharp}
private void BtnToggle_Click(object? sender, EventArgs e)
{
    if (dgvDevices.SelectedRows.Count == 0)
    {
        MessageBox.Show("Выберите устройство", "Внимание",
            MessageBoxButtons.OK, MessageBoxIcon.Warning);
        return;
    }

    int deviceId = Convert.ToInt32(
        dgvDevices.SelectedRows[0].Cells["Id"].Value);
    var deviceRow = DatabaseHelper.GetDeviceById(deviceId);

    if (deviceRow == null) return;

    bool currentState = Convert.ToInt32(
        deviceRow["IsOn"]) == 1;
    DatabaseHelper.SetDeviceState(deviceId, !currentState);
    LoadDevices();
}
\end{minted}

\section*{Алгоритм: Включение/выключение устройства}

\begin{enumerate}
    \item Пользователь выбирает устройство в списке \texttt{DataGridView}.
    \item Нажимает кнопку <<Вкл/Выкл>>.
    \item Проверить, что строка выбрана.
    \item Взять Id из выбранной строки.
    \item Подключиться к SQLite через \texttt{DatabaseHelper}.
    \item Получить текущее состояние устройства:
    \[
    \text{SELECT IsOn FROM Devices WHERE Id = ?}
    \]
    \item Инвертировать состояние и выполнить UPDATE:
    \[
    \text{UPDATE Devices SET IsOn = ? WHERE Id = ?}
    \]
    \item Перезагрузить таблицу на экране.
\newpage
    \begin{table}[h!]
\centering
\begin{tabular}{|l|p{10cm}|}
\hline
\textbf{Функция} & \textbf{Описание} \\
\hline
GetAllDevices() & SELECT * FROM Devices, возвращает DataTable \\
\hline
GetDeviceById(id) & SELECT с WHERE Id = ?, возвращает словарь \\
\hline
InsertDevice(device) & INSERT в Devices, возвращает сгенерированный Id \\
\hline
DeleteDevice(id) & DELETE из Devices и дочерней таблицы \\
\hline
UpdateDevice(device, id) & UPDATE Name, Manufacturer, IsOn \\
\hline
SetDeviceState(id, isOn) & UPDATE IsOn для быстрого вкл/выкл \\
\hline
CreateDeviceObject(id) & Создаёт объект C# нужного типа по данным из БД \\
\hline
SaveDeviceObject(device, id) & Сохраняет изменения объекта в БД \\
\hline
\end{tabular}
\caption*{}
\label{tab:db_methods}
\end{table}
\end{enumerate}

\textbf{GetAllDevices()}

\begin{verbatim}
// MainForm.cs — LoadDevices()
var dt = DatabaseHelper.GetAllDevices();
foreach (DataRow row in dt.Rows)
    dgvDevices.Rows.Add(row["Id"], row["Name"],
        row["DeviceType"], row["Manufacturer"], ...);
\end{verbatim}

\textbf{GetDeviceById()}

\begin{verbatim}
// MainForm.cs — ShowDeviceDetails()
var deviceRow = DatabaseHelper.GetDeviceById(deviceId);
string name = deviceRow["Name"].ToString();
string type = deviceRow["DeviceType"].ToString();
\end{verbatim}

\textbf{InsertDevice()}

\begin{verbatim}
// DeviceEditForm.cs — BtnSave_Click()
int deviceId = DatabaseHelper.InsertDevice(device);

switch (type)
{
    case "Thermostat":
        DatabaseHelper.InsertThermostat(
            (Thermostat)device, deviceId);
        break;
    case "Light":
        DatabaseHelper.InsertLight(
            (Light)device, deviceId);
        break;
    // ...аналогично для DoorLock, Camera
}
\end{verbatim}

\textbf{DeleteDevice()}

\begin{verbatim}
// MainForm.cs — BtnDelete_Click()
int deviceId = Convert.ToInt32(
    dgvDevices.SelectedRows[0].Cells["Id"].Value);
DatabaseHelper.DeleteDevice(deviceId);
// DELETE FROM <тип> WHERE DeviceId = ?
// DELETE FROM Devices WHERE Id = ?
\end{verbatim}

\textbf{UpdateDevice() / SaveDeviceObject()}

\begin{verbatim}
// DeviceEditForm.cs — BtnSave_Click() (режим редактирования)
DatabaseHelper.UpdateDevice(device, _deviceId);

// MethodCallForm.cs — BtnExecute_Click()
DatabaseHelper.SaveDeviceObject(device, _deviceId);
// Обновляет Devices + дочернюю таблицу
\end{verbatim}

\textbf{CreateDeviceObject()}

\begin{verbatim}
// MethodCallForm.cs — создание объекта для вызова методов
var device = DatabaseHelper.CreateDeviceObject(_deviceId);
// Создаёт Thermostat/Light/DoorLock/Camera

((Thermostat)device).SetTemperature(temp);
DatabaseHelper.SaveDeviceObject(device, _deviceId);
\end{verbatim}

\textbf{LoadDevices()}
\begin{table}[h!]
\centering
\begin{tabular}{|c|l|}
\hline
\textbf{Шаг} & \textbf{Действие} \\
\hline
1 & Очищает строки DataGridView \\
\hline
2 & Инициализирует DatabaseHelper (если не инициализирован) \\
\hline
3 & Выполняет GetAllDevices() — SELECT * FROM Devices \\
\hline
4 & Перебирает строки DataTable, добавляет в DataGridView \\
\hline
5 & Отображает: ID, Название, Тип, Производитель, Статус \\
\hline
\end{tabular}
\caption*{}
\end{table}


\newpage
\textbf{DeviceEditForm.cs}\\
\begin{enumerate}
\item DeviceEditForm() — Конструктор для создания нового устройства\\
\item DeviceEditForm(int deviceId) — Конструктор для редактирования существующего
\item InitializeComponent() — Создание UI: выбор типа (ComboBox), выбор конструктора (простой/расширенный), поля для параметров каждого типа
\item CmbType\_SelectedIndexChanged() — Переключение панели с полями при смене типа устройства
\item LoadDeviceData() — Загрузка данных существующего устройства в форму
\item BtnSaveClick() — Сохранение: создание объекта C# нужного типа, INSERT или UPDATE в БД
\end{enumerate}\\

\textbf{MethodCallForm.cs}\\
\begin{enumerate}
\item MethodCallForm(int deviceId) — Конструктор, принимает Id устройства
\item InitializeComponent() — Создание UI: ComboBox со списком методов, поле для параметра, кнопка "Выполнить"
\item LoadMethods() — Загружает список методов в зависимости от типа устройства (Thermostat → SetTemperature/GetClimateInfo, Light → SetBrightness/GetLightInfo, и т.д.)
\item CmbMethods\_SelectedIndexChanged() — Показ/скрытие поля параметра
\item BtnExecuteClick() — Вызов метода на объекте из CreateDeviceObject(), сохранение через SaveDeviceObject()
\end{enumerate}
\newpage
\textbf{MainForm.cs}\\

\textbf{Методы:}

\begin{enumerate}
\item MainForm() — Конструктор: инициализация UI + загрузка устройств\\
\item InitializeComponent() — Создание UI: DataGridView, панель кнопок, панель деталей\\
\item LoadDevices() — Загрузка списка устройств из БД в DataGridView\\
\item ShowDeviceDetails(int deviceId) — Детальная информация о выбранном устройстве
\end{enumerate}\\


\textbf{UI обработчики}\\

\begin{enumerate}
\item BtnAddClick() — Открыть DeviceEditForm для добавления\\
\item BtnEditClick() — Открыть DeviceEditForm для редактирования\\
\item BtnDeleteClick() — Удалить выбранное устройство\\
\item BtnToggleClick() — Включить/выключить выбранное устройство\\
\item BtnMethodsClick() — Открыть MethodCallForm для вызова методов\\
\item BtnRefreshClick() — Обновить список устройств\\
\item DgvDevicesSelectionChanged() — Обновление панели деталей при выборе
\end{enumerate}\\

\begin{figure}[H]
    \centering
    \includegraphics[width=0.9\linewidth]{image2.png}
    \caption*{Рисунок 3 — ERD диаграмма}
    \label{fig:placeholder}
\end{figure}
\newpage
\begin{figure}[H]
    \centering
    \includegraphics[width=1\linewidth]{image3.png}
    \caption*{Рисунок 4 — Экранные формы}
    \label{fig:placeholder}
\end{figure}

\begin{figure}[H]
    \centering
    \includegraphics[width=1\linewidth]{image4.png}
    \caption*{Рисунок 5 — Экранные формы}
    \label{fig:placeholder}
\end{figure}

\begin{figure}[H]
    \centering
    \includegraphics[width=1\linewidth]{image5.png}
    \caption*{Рисунок 6 — Экранные формы}
    \label{fig:placeholder}
\end{figure}

\begin{figure}[H]
    \centering
    \includegraphics[width=1\linewidth]{image6.png}
    \caption{Рисунок 7 — Экранные формы}
    \label{fig:placeholder}
\end{figure}

\begin{figure}[H]
    \centering
    \includegraphics[width=1\linewidth]{image7.png}
    \caption*{Рисунок 8 — Экранные формы}
    \label{fig:placeholder}
\end{figure}

\section*
{Вывод}
В рамках лабораторной работы была спроектирована иерархия классов, включающая один абстрактный базовый класс \texttt{Device} и четыре производных от него: \texttt{Thermostat}, \texttt{Light}, \texttt{DoorLock}, \texttt{Camera}. Реализовано приложение на Windows Forms с серверной логикой на C\#. С использованием библиотеки System.Data.SQLite создана база данных из пяти таблиц: \texttt{Devices}, \texttt{Thermostats}, \texttt{Lights}, \texttt{DoorLocks}, \texttt{Cameras}. Данные хранятся в отдельных таблицах для каждого типа устройств, связанных через внешний ключ \texttt{DeviceId}. Реализованы все основные CRUD-функции, а также методы классов, доступные через графический интерфейс (включение/выключение, установка температуры, управление яркостью, открытие замка, запуск записи и другие).

\end{document}
